\enlargethispage{2\baselineskip}%
% Resumo (NBR 6028: um parágrafo, 150-500 palavras, sem fórmulas).
% Reescrito em 24/11/2026 (revisão pedagógica); todos os números conferidos
% nos capítulos e artefatos de origem antes da inclusão.
Por uma década, a escolha de modelos para dados tabulares apoiou-se na
máxima de que o \emph{gradient boosting} sempre vence. Os
\emph{foundation models} tabulares, redes pré-treinadas que predizem sem
treinamento na tarefa, puseram-na em dúvida. Contudo, a evidência vem sobretudo dos proponentes, sob protocolos
distintos. Esta dissertação produz evidência sob protocolo único e a
converte em procedimento de escolha, em três frentes. Primeiro, um
benchmark multicritério avalia 14 modelos de três famílias
(\emph{gradient boosting}, aprendizado profundo e \emph{foundation
models}) em 18 conjuntos do OpenML, com \emph{hold-out} intocado pela
seleção, inferência frequentista e bayesiana e custos de treino, ajuste e
latência medidos. Segundo, avalia-se um arcabouço de decisão: o componente
aprendido, sob validação \emph{leave-one-dataset-out}; as políticas fixas,
descritivamente. Terceiro, estuda-se a destilação das curvas de quantis do
professor TabPFN em alunos de \emph{gradient boosting}, em 20 conjuntos de
regressão, sob hipóteses registradas. Entre os 11 modelos do núcleo
(anteriores à extensão de 2026), o topo da acurácia na
classificação binária é um empate. As demais tarefas são lidas
descritivamente. Sob uma ROPE de um ponto percentual de AUC, 3 dos outros
10 modelos são praticamente equivalentes ao líder, o TabPFN, ele próprio um \emph{foundation model}, pelo critério bayesiano estrito, e 8 pelo brando, afirmação bem mais fraca. Na extensão de 2026, o TabFM, sem ajuste, tem o melhor
rank médio nas três tarefas (na multiclasse, sobre os dois conjuntos em que
opera) e é o melhor modelo em 13 dos 18 conjuntos (opera em 17). A
liderança é consistente entre conjuntos, mas os testes, de baixo poder, não
a confirmam como separação estatística, e ela custa cerca de cinco ordens de magnitude a mais em latência que o XGBoost no conjunto de referência (\texttt{adult}). As
redes Kolmogorov--Arnold ficam no último terço nas três tarefas. Na
decisão, um roteador aprendido por meta-características acerta menos que a
regra trivial (0,11 contra 0,44). A política fixa ``\emph{foundation
model} primeiro'' (um oráculo intrafamília) tem \emph{regret} mediano zero
na métrica primária, mas deixa de ser a recomendada quando 10 a 20
por cento do peso passa da acurácia à latência, pelo \emph{regret} mediano, na normalização por postos. Na destilação, o
aluno treinado nas médias do professor supera o controle em só 4 de 20. O aluno treinado nas curvas de quantis retém, na mediana, 12 a
19 por cento da vantagem distribucional, conforme o ajuste, com retenção positiva em 16 de 20 (teste de sinal, p aproximadamente 0,019
após a correção sequencial). Ele supera \emph{baselines} distribucionais
clássicos, em configuração \emph{default}, em 17 de 20 (12 de 20 contra a melhor alternativa por conjunto, via oráculo). No regime estudado (até 100 mil amostras,
versões fixadas, GPU única), a escolha deve partir não de um
vencedor universal, mas das restrições da aplicação: latência, limites
arquiteturais e explicações auditáveis. Nesse regime, a destilação
recupera, na mediana, uma parcela modesta da incerteza preditiva do
\emph{foundation model}, sob pré-condições mensuráveis antes do treino do
aluno. A regra de decisão foi avaliada apenas nos conjuntos de que foi derivada, e sua pré-condição, só em versão oráculo.

\vspace{0.4cm}
\noindent\textbf{Palavras-chave:} dados tabulares; foundation models;
benchmark; destilação de conhecimento; regressão distribucional.
